% Représentations d'une molécule et isomères de squelette (L'essentiel, chapitre 11)
\begin{tikzpicture}[x=1cm,y=1cm]
\node[anchor=west] at (0,1.4) {\chemfig[atom sep=1.9em]{CH_3-CH(-[::60]OH)-CH_2-CH_3}};
\draw[pcFleche] (3.65,1.4) -- (4.55,1.4);
\node[anchor=west] at (4.8,1.3) {\chemfig[atom sep=1.9em]{-[:30](-[:90]OH)-[:-30]-[:30]}};
\node[pcLegende,anchor=north,align=center] at (3.4,0.55) {butan-2-ol : formule semi-développée et formule topologique};
\node[anchor=west] at (8.3,1.75) {\chemfig[atom sep=1.9em]{-[:30]-[:-30]-[:30]}};
\node[anchor=west] at (8.45,0.75) {\chemfig[atom sep=1.9em]{-[:30](-[:90])-[:-30]}};
\node[pcLegende,anchor=west,align=left] at (11.0,1.85) {butane};
\node[pcLegende,anchor=west,align=left] at (11.0,0.95) {2-méthylpropane};
\node[pcLegende,anchor=north,align=center] at (10.6,0.25) {deux squelettes pour \ce{C4H10} :\\des isomères de constitution};
\end{tikzpicture}
\par\smallskip\legende{Dans une formule topologique, chaque extrémité ou brisure de ligne représente un atome de carbone portant ses hydrogènes.}
